\section*{Chapitre \ref{ch:iv2:portfolio-construction} --- Construction de portefeuille}
\markboth{\protect\omTocMark{}{Réponses~: Construction de portefeuille}}{\protect\omTocMark{}{Réponses~: Construction de portefeuille}}

\begin{solution}{iq:iv2:portfolio-construction:1}
Les poids valent environ 85.7~\% et 14.3~\%, et la volatilité vaut environ 19.6~\%.

$w = (\sigma_2^2 - \sigma_{12})/(\sigma_1^2 + \sigma_2^2 - 2\sigma_{12})$, avec $\sigma_{12} = 0.03$. Donc $w = 0.06/0.07 = 6/7$.

La variance vaut $w^2(0.04) + (1 - w)^2(0.09) + 2w(1 - w)(0.03) = 0.0386$.

\iqwrong{Allouer selon l'inverse de la variance seule. Cela ignore la corrélation et donne 69/31, qui a ici une variance plus élevée.}
\end{solution}

\begin{solution}{iq:iv2:portfolio-construction:2}
C'est 50/50, avec une volatilité d'environ 15.5~\%.

Par symétrie, les poids sont égaux, et la variance vaut $\sigma^2(1 + \rho)/2 = 0.04 \times 0.6 = 0.024$. Sa racine vaut 15.49~\%.

\iqwrong{Dire que la corrélation modifie les poids. À volatilités égales, elle ne modifie que le risque résultant.}
\end{solution}

\begin{solution}{iq:iv2:portfolio-construction:3}
Les poids valent 2/3 et 1/3, et non, la corrélation n'importe pas pour deux actifs.

Une contribution vaut $w_i(\Sigma w)_i$. Pour deux actifs, en égalisant les deux, on obtient $w_1^2\sigma_1^2 + \sigma_{12}w_1w_2 = w_2^2\sigma_2^2 + \sigma_{12}w_1w_2$. Le terme croisé s'annule, donc $w_1\sigma_1 = w_2\sigma_2$.

\iqwrong{Transposer cela à trois actifs ou plus. Là, l'inverse de la volatilité ne donne un risque égal que si toutes les corrélations sont égales.}
\end{solution}

\begin{solution}{iq:iv2:portfolio-construction:4}
Il vaut 0.50 avec 100 paris, et environ 0.79 avec 252.

$\mathrm{IR} = \mathrm{IC}\sqrt{\mathrm{BR}}$~: $0.05 \times 10 = 0.5$, et $0.05 \times 15.87 = 0.79$.

\iqwrong{Compter les transactions comme des paris. La largeur compte les prévisions indépendantes par an, et une action détenue sur 252 jours à partir d'un seul signal lent ne constitue pas 252 paris.}
\end{solution}

\begin{solution}{iq:iv2:portfolio-construction:5}
Le turnover à sens unique vaut 10~\%, et il coûte 2 points de base, soit environ 24 points de base par an.

La somme des variations absolues vaut 0.2, et sa moitié, soit le turnover à sens unique, vaut 0.1. On échange 0.2 dollar, on vend 0.1 et on achète 0.1, à 10 points de base~: $0.2 \times 10 = 2$ points de base de portefeuille. Douze fois donne 24.

\iqwrong{Chiffrer le nombre à sens unique. On paie à la fois sur la vente et sur l'achat, et les cotations de turnover varient, donc il faut vérifier celle qui est visée.}
\end{solution}

\begin{solution}{iq:iv2:portfolio-construction:6}
On vend à découvert 15~millions de dollars, et la position nette en dollars est de $-\$5$ millions.

La position longue a des dollars bêta de $10 \times 1.2 = 12$. Donc $12/0.8 = 15$.

\iqwrong{Supposer que neutre en bêta et neutre en dollars sont la même chose. Ici, être neutre sur l'un revient à être à découvert de 5~millions de dollars sur l'autre. Si l'estimation du bêta est fausse, la couverture l'est aussi.}
\end{solution}

\begin{solution}{iq:iv2:portfolio-construction:7}
Les poids de tangence valent environ 72.4~\% et 27.6~\%, et le ratio de Sharpe vaut environ 0.429. Avec une aversion au risque de 4, ils détiennent 38.9~\% et 14.8~\%, avec un rendement d'environ 4.59~\% et une volatilité d'environ 10.7~\%.

On résout $\Sigma x = \mu$. Ici $x = (1.556, 0.593)$. La normalisation donne les poids de tangence. Le ratio de Sharpe vaut $\sqrt{\mu^\top x} = \sqrt{0.1837} = 0.429$.

L'investisseur détient $x/\lambda$, soit $(0.389, 0.148)$, et le reste est en cash. Donc le rendement vaut $0.1837/4 = 4.59~\%$, et la volatilité vaut $0.429/4 = 10.7~\%$.

\iqwrong{Forcer les poids à totaliser 1 pour l'investisseur. Avec un actif sans risque, l'aversion au risque fixe la taille~: elle met à l'échelle tout le portefeuille risqué, et le reste est en cash.}
\end{solution}

\begin{solution}{iq:iv2:portfolio-construction:8}
Les poids valent environ 123.5~\% et $-23.5~\%$, avec une volatilité d'environ 19.5~\%. Le second poids est négatif lorsque $\rho > \sigma_1/\sigma_2$, et ici $0.8 > 0.667$.

La covariance vaut $0.048$, donc $w = (0.09 - 0.048)/(0.13 - 0.096) = 1.235$.

Le second poids vaut $(\sigma_1^2 - \rho\sigma_1\sigma_2)$ divisé par un dénominateur positif. Il est négatif lorsque $\sigma_1 < \rho\sigma_2$. À forte corrélation, l'actif le plus risqué est surtout une copie à effet de levier de l'autre. Une vente à découvert de celui-ci couvre donc le premier.

\iqwrong{Traiter une vente à découvert dans un portefeuille de variance minimale comme une erreur. C'est la bonne couverture, mais elle repose sur une estimation de la corrélation qui peut ne pas tenir en période de stress.}
\end{solution}

\begin{solution}{iq:iv2:portfolio-construction:9}
Le ratio de Sharpe vaut environ 0.674, avec $x = (5.91, -2.73)$. Le meilleur ratio de Sharpe en position longue uniquement vaut 0.50, entièrement en A.

On résout $\Sigma x = \mu$, avec des variances de 0.04 et 0.09 et une covariance de 0.05. Donc $x = (5.91, -2.73)$, et le ratio de Sharpe vaut $\sqrt{\mu^\top x} = 0.674$.

En position longue uniquement, la solution sans contrainte veut un B négatif. Le ratio de Sharpe est unimodal le long du mélange de A et B, donc le meilleur en position longue se trouve à la frontière, sans B. Et $0.10/0.20 = 0.5$.

Pouvoir vendre B à découvert ajouterait donc 0.17 de ratio de Sharpe, mais à un levier élevé~: 5.91 en position longue et 2.73 en position courte par unité d'aversion au risque, sur une covariance qui doit être juste.

\iqwrong{Affirmer qu'un actif à rendement élevé doit figurer dans tout bon portefeuille. B a un rendement positif et il est tout de même vendu à découvert ici, car il est fortement corrélé à A et a un moins bon ratio de Sharpe.}
\end{solution}

\begin{solution}{iq:iv2:portfolio-construction:10}
Elle vaut 18.2~\%, 17.1~\% et 16.6~\%, et la limite est de 16.4~\%. Il faut environ 19 actions pour rester à un point près.

La variance vaut $\sigma^2(1/n + (1 - 1/n)\rho)$. Lorsque $n \to \infty$, $\sigma\sqrt{\rho} = 30~\% \times 0.548 = 16.43~\%$.

À un point près signifie en dessous de 17.43~\%. Donc $0.09(0.7/n + 0.3) \le 0.1743^2 = 0.03038$, et $n \ge 19$.

\iqwrong{Dire qu'un nombre plus élevé d'actions signifie toujours beaucoup moins de risque. Les 19 premières font presque tout le travail, et le facteur commun demeure, quel que soit le nombre ajouté.}
\end{solution}

\begin{solution}{iq:iv2:portfolio-construction:11}
La largeur effective vaut environ 83, et le ratio d'information environ 0.18. Si on le traite comme indépendant, il vaut environ 0.45.

Avec $N$ paris de corrélation deux à deux $\rho$, la variance de la moyenne vaut $(1 + (N - 1)\rho)/N$, donc le nombre effectif vaut $N/(1 + (N - 1)\rho) = 500/5.99 = 83.5$. Donc $\mathrm{IR} = 0.02\sqrt{83.5} = 0.183$.

Même une corrélation de 0.01 divise la largeur par six. Un pari sur un facteur déguisé en 500 titres présente ce problème.

\iqwrong{Compter les actions pour la largeur. Une petite erreur commune sur de nombreux titres réduit fortement ce nombre, et à la limite donne $1/\rho$ paris indépendants, soit ici 100.}
\end{solution}

\begin{solution}{iq:iv2:portfolio-construction:12}
Il vaut environ 0.34.

$\mathrm{IR} = \mathrm{TC} \times \mathrm{IC}\sqrt{\mathrm{BR}} = 0.3 \times 0.05 \times 22.36 = 0.335$.

Le coefficient de transfert est la corrélation entre les poids actifs que vous détenez et les poids que votre prévision seule donnerait. Il mesure la part du signal que les contraintes laissent passer.

Pour l'augmenter, il faut relâcher la contrainte long-only (un portefeuille long-only ne peut sous-pondérer un petit titre que de son poids dans l'indice de référence), élargir l'univers ou réduire les limites de tracking error. Une poche long-short peut l'augmenter le plus.

\iqwrong{Améliorer d'abord l'IC. Avant de chercher plus de signal, il faut vérifier combien on perd à la construction~: ici 70~\%.}
\end{solution}

\begin{solution}{iq:iv2:portfolio-construction:13}
L'erreur type de chaque moyenne vaut environ 4.5~\%. Avec le shrinkage, on obtient 10~\% et 0~\%. L'optimisation moyenne--variance donne alors des paris beaucoup moins extrêmes, car elle amplifie moins les erreurs.

$20~\%/\sqrt{20} = 4.47~\%$. Une moyenne de 15~\% a un $t$ de 3.4, et l'écart entre les deux a une erreur type de 6.3~\%, soit un $t$ de 3.2.

Le shrinkage divise par deux l'écart, qui passe de 20 à 10 points. Les poids sont linéaires en les moyennes et ne dépendent que de leurs différences à risque égal. Donc l'inclinaison long-short est divisée par deux.

\textbf{Pourquoi il maximise les erreurs.} L'optimiseur surpondère les actifs à rendement estimé élevé et risque estimé faible. Ce sont probablement ceux où l'erreur d'estimation joue le plus en faveur du portefeuille. Ainsi la performance dans l'échantillon (in-sample) est surestimée, et la performance hors échantillon (out-of-sample) déçoit.

\iqwrong{Penser que vingt ans, c'est beaucoup de données. La moyenne dépend de la durée, et non du nombre d'observations. Les données à haute fréquence aident beaucoup pour la variance, mais presque pas pour la moyenne.}
\end{solution}

\begin{solution}{iq:iv2:portfolio-construction:14}
Le conditionnement vaut 500. Après shrinkage, les valeurs propres valent 4.40, 1.20 et 0.41, et le conditionnement vaut environ 10.8.

La trace moyenne vaut $6.01/3 = 2.003$. Chaque valeur propre vaut donc $0.8\lambda + 0.2 \times 2.003$, les vecteurs propres étant inchangés.

Un portefeuille de variance minimale est $\propto \Sigma^{-1}\mathbf{1}$, et l'inverse amplifie les petites valeurs propres. Une valeur propre de 0.01 est en général du bruit~: elle signale une combinaison quasi sans risque, dans laquelle l'optimiseur s'engouffre. La relever à 0.41 divise ce poids par un facteur 41.

\iqwrong{Traiter une petite valeur propre comme réelle. Son estimation sur échantillon est biaisée à la baisse quand les actifs sont nombreux par rapport aux données, et avec plus d'actifs que de jours, certaines sont exactement nulles.}
\iqfollow{Comment choisissez-vous l'intensité du shrinkage~?}{En minimisant l'erreur quadratique moyenne attendue, dont on dispose d'une estimation en forme close. Ou en testant la variance hors échantillon du portefeuille optimisé.}
\end{solution}

\begin{solution}{iq:iv2:portfolio-construction:15}
Avec 100 titres et 1\,000 jours, les valeurs propres se situent à peu près dans $[0.47, 1.73]$. Avec 500, la borne haute monte à environ 2.91, et la borne basse à environ 0.09.

Avec $q = N/T$, les bornes valent $(1 \pm \sqrt{q})^2$. Pour $q = 0.1$, cela donne $0.468$ et $1.732$. Pour $q = 0.5$, on obtient $0.086$ et $2.914$.

Une matrice de bruit pur peut encore avoir une valeur propre de 1.7 ou de 2.9, et même une vraie identité paraît structurée. Les données réelles présentent une grande valeur propre de marché bien au-delà de la bande, et peut-être quelques valeurs propres sectorielles.

On l'utilise comme filtre~: les valeurs propres situées dans la bande de bruit sont traitées comme du bruit et remplacées par leur moyenne, ce qui conserve la trace.

\iqwrong{Utiliser chaque composante principale d'un échantillon comme facteur. Celles qui sont dans la bande ne portent guère plus que du bruit, et elles rendent un backtest meilleur qu'il ne l'est.}
\iqfollow{Quelle hypothèse ce résultat peut-il remettre en cause~?}{Les bandes supposent une indépendance dans le temps et des queues fines. Les queues épaisses et le regroupement de volatilité élargissent le spectre de l'échantillon, si bien que le vrai seuil est plus élevé.}
\end{solution}

\begin{solution}{iq:iv2:portfolio-construction:16}
Il y en a 125\,250 pour la matrice complète, et 5\,555 pour le modèle factoriel.

Complète~: $500 \times 501/2$. Factoriel~: $500 \times 10$ charges, $10 \times 11/2 = 55$ covariances factorielles et 500 variances idiosyncratiques.

L'hypothèse est que les résidus sont non corrélés une fois les facteurs retirés. Cela donne $\Sigma = BFB^\top + D$. Elle rend aussi la matrice inversible, même avec moins de jours que de titres.

Elle échoue lorsque les résidus sont en fait liés~: relations de paires, facteur commun manquant comme une position encombrée (crowded trade), ou une fusion. Il faut donc vérifier les corrélations résiduelles, et ne pas considérer un faible risque factoriel comme un faible risque lorsqu'un book est vendeur à découvert d'un résidu commun.

\iqwrong{Dire que les modèles factoriels ne sont qu'une approximation faute de données. Ils régularisent aussi le problème, si bien qu'ils battent souvent la matrice empirique hors échantillon, même avec beaucoup de données.}
\end{solution}

\begin{solution}{iq:iv2:portfolio-construction:17}
Chaque actif a une croissance logarithmique espérée de $-2~\%$ par an, et le portefeuille rééquilibré environ $-1~\%$. La stratégie d'achat et conservation obtient environ $-1.2~\%$ sur 30 ans. Avec une corrélation de 1, le rééquilibrage ne rapporte rien.

La croissance logarithmique est le rendement arithmétique diminué de la moitié de la variance. Chaque actif a $0 - 0.04/2 = -2~\%$. Le portefeuille a une variance de $0.02$, donc $0 - 0.01 = -1~\%$.

Le gain associé à chaque actif pris isolément est donc $\tfrac12(\overline{\sigma_i^2} - \sigma_p^2) = 1\%$, soit le \index{rendement de diversification}rendement de diversification. C'est ce qu'on obtient en comparant un portefeuille équilibré à l'actif moyen.

Un portefeuille d'achat et conservation part lui aussi diversifié, mais il dérive vers le gagnant et finit par ressembler à un seul actif. Ainsi, sur 30 ans, le taux de croissance logarithmique espéré vaut environ $-1.2~\%$, et l'écart avec le rééquilibrage n'est que d'environ 0.2~\% par an dans ce cadre. Les coûts peuvent effacer ce gain.

Avec une corrélation de 1, les deux actifs sont identiques, la variance du portefeuille ne change pas, et aucun gain n'apparaît.

\iqwrong{Dire que le rééquilibrage est de l'argent gratuit~: il relève la croissance logarithmique en réduisant la pénalité de variance, et non en ajoutant du rendement. La richesse finale espérée n'est pas plus élevée ici, et elle est plus faible lorsqu'un actif suit une tendance.}
\end{solution}

\begin{solution}{iq:iv2:portfolio-construction:18}
Le ratio de Sharpe vrai vaut environ 0.57. Le portefeuille tangent estimé affiche environ 1.65 dans l'échantillon, et seulement environ 0.19 hors échantillon. On a donc perdu les deux tiers de la vérité, tout en affichant trois fois la vraie valeur.

\begin{sloppypar}\textbf{Vérité.} Avec des poids égaux, le rendement vaut 0.5~\% par mois, et la volatilité vaut $5~\%\sqrt{0.1 + 0.9 \times 0.3} = 3.04~\%$. Le ratio de Sharpe mensuel vaut donc 0.164, soit 0.569 par an. Par symétrie, c'est le véritable optimum.\end{sloppypar}

\textbf{Estimé.} Avec 10 actifs et 60 mois, l'erreur sur les moyennes est du même ordre de grandeur que les moyennes elles-mêmes. Un ratio de Sharpe empirique est élevé parce que les poids exploitent le bruit. Dans une simulation de 3\,000 répétitions, le ratio de Sharpe dans l'échantillon vaut en moyenne 1.65, et les mêmes poids évalués sur les vrais paramètres valent en moyenne 0.19.

Ainsi~: les poids égaux ou un shrinkage fort battent l'optimisation sur échantillon, à moins que le signal ne soit beaucoup plus fort que cela. Il faut beaucoup plus d'observations par actif, ou une structure sur les moyennes.

\iqwrong{Accuser la seule covariance. L'essentiel du dommage vient ici des moyennes. Si l'on savait que les différences de moyennes étaient nulles, la covariance estimée permettrait encore de se rapprocher beaucoup de l'optimum.}
\end{solution}

\begin{solution}{iq:iv2:portfolio-construction:19}
Les rendements implicites valent 9.0~\% et 13.5~\%. Les rendements a posteriori valent 9.5~\% et 13.875~\%, ce qui donne des poids de 65~\% et 40~\%. Lorsque l'incertitude tend vers zéro, les rendements a posteriori tendent vers 10~\% et 14.25~\%, et les poids vers 70~\% et 40~\%.

\textbf{A priori.} $\Pi = \delta\Sigma w = 2.5\,(0.036, 0.054) = (0.09, 0.135)$.

\textbf{A posteriori.} La mise à jour s'écrit
\[
\mu = \Pi + \tau\Sigma P^\top(P\tau\Sigma P^\top + \Omega)^{-1}(Q - P\Pi).
\]
Ici $P = (1, 0)$, l'écart vaut $0.01$, $\tau\Sigma_{11} = 0.002$, et $\Omega = 0.002$. Ainsi, l'actif 1 bouge de la moitié de l'écart, soit $+0.005$. L'actif 2 bouge de $(0.0015/0.004) \times 0.01 = 0.00375$, par l'intermédiaire de la corrélation.

\textbf{Poids.} En divisant par $\delta\Sigma$, l'actif 2 reste à 0.4, et seul l'actif 1 bouge. C'est une propriété connue~: seuls les actifs faisant l'objet d'une vue s'écartent du poids de marché. Lorsque $\Omega \to 0$, l'actif 1 atteint entièrement sa vue, jusqu'au poids 0.70. Lorsque $\Omega \to \infty$, on retrouve l'a priori et le portefeuille de marché.

\iqwrong{Dire que le rendement du deuxième actif ne change pas parce qu'on n'a aucune vue sur lui. Son rendement bouge par la corrélation, et c'est pourquoi son poids reste le même.}
\end{solution}

\begin{solution}{iq:iv2:portfolio-construction:20}
Les poids égaux donnent des parts de 29.8~\%, 14.3~\% et 56.0~\%. La pondération inverse à la volatilité donne 35.2~\%, 38.3~\% et 26.6~\%. Le troisième ensemble donne un tiers à chaque actif, ce qui en fait le portefeuille à risque égal, avec une volatilité d'environ 11.6~\%.

Une part de risque vaut $w_i(\Sigma w)_i/w^\top\Sigma w$. Pour le dernier ensemble, $\Sigma w = (0.0162, 0.0087, 0.0211)$ à quatre décimales près. Les contributions au risque valent donc toutes 0.00447, et la somme des trois est la variance, soit 0.0134.

La pondération inverse à la volatilité ignore la corrélation. Elle ne serait juste ici que si les corrélations étaient égales. Le deuxième actif a une faible volatilité et est positivement corrélé aux deux autres, si bien qu'il reçoit un poids trop élevé selon cette règle.

La parité de risque utilise la corrélation, et elle n'utilise pas les rendements espérés. C'est un véritable argument~: elle est moins sensible à des entrées de mauvaise qualité que l'optimisation moyenne--variance. Elle reste néanmoins dépendante de la covariance, et elle a tendance à être employée avec un effet de levier pour atteindre un objectif de rendement.

\iqwrong{Penser que des contributions au risque égales signifient une volatilité égale des positions. Il s'agit des contributions au total, une fois la corrélation prise en compte.}
\iqfollow{Existe-t-il toujours une solution~?}{Oui, pour un portefeuille long-only à variance positive. On minimise $\tfrac12 w^\top\Sigma w - c\sum_i\ln w_i$, puis on normalise pour que la somme vaille 1~: c'est un problème convexe qui admet une seule réponse.}
\end{solution}

\begin{solution}{iq:iv2:portfolio-construction:21}
Le portefeuille de variance minimale globale a un rendement d'environ 10.2~\% et une volatilité d'environ 14.0~\%. Pour un rendement de 12~\%, les poids valent environ $(0.019, -0.029, 1.010)$, et la volatilité environ 15.0~\%.

Posons $A = \mathbf{1}^\top\Sigma^{-1}\mathbf{1} = 51.29$, $B = \mathbf{1}^\top\Sigma^{-1}\mu = 5.243$, $C = \mu^\top\Sigma^{-1}\mu = 0.6437$, et $D = AC - B^2 = 5.531$.

Le minimum est atteint pour un rendement $B/A = 10.22~\%$ et une volatilité $1/\sqrt{A} = 13.96~\%$.

Pour un objectif $m$, $w = \Sigma^{-1}(\lambda\mathbf{1} + \gamma\mu)$ avec $\lambda = (C - Bm)/D$ et $\gamma = (Am - B)/D$. La variance vaut $(Am^2 - 2Bm + C)/D$.

\textbf{Deux fonds.} Tout portefeuille de la frontière est un mélange de $\Sigma^{-1}\mathbf{1}$ et de $\Sigma^{-1}\mu$, avec des coefficients linéaires en $m$. Deux portefeuilles quelconques de la frontière engendrent donc celle-ci, et il n'est pas nécessaire de recalculer la solution pour chaque objectif.

Un rendement de 10~\% est inférieur au rendement du portefeuille de variance minimale~: ce point se trouve donc sur la partie inefficiente, car un rendement plus élevé est accessible au même risque.

\iqwrong{Croire qu'un objectif plus bas signifie moins de risque partout. Sous le rendement de la variance minimale, le risque remonte.}
\end{solution}

\begin{solution}{iq:iv2:portfolio-construction:22}
Il vaut $(0.4, 0.2, 0.4)$, avec une volatilité d'environ 14.9~\%.

Le troisième poids étant plafonné à 0.4, les autres augmentent, et le premier atteint lui aussi le plafond~: le second reçoit le reste.

\textbf{Conditions.} Avec $w^\top\mathbf{1} = 1$ et $w_i \le 0.4$, les conditions s'écrivent $\Sigma w = \lambda\mathbf{1} - \nu$, avec $\nu_i \ge 0$, et $\nu_i = 0$ sauf si $w_i$ est au plafond. Ici $\Sigma w = (0.0232, 0.0300, 0.0174)$ à quatre décimales près. L'actif situé sous le plafond est le second, donc $\lambda = 0.0300$. Les autres ont $\nu_1 = 0.0068$ et $\nu_3 = 0.0126$. Ces deux valeurs étant positives, les conditions sont satisfaites.

Le problème étant convexe, ces conditions sont suffisantes~: c'est l'unique solution optimale.

Il faut lire les multiplicateurs. Ce sont les gains obtenus en relâchant chaque plafond. Relâcher le plafond de l'actif 3 rapporte presque deux fois plus que relâcher celui de l'actif 1.

\iqwrong{Plafonner le poids élevé et redimensionner les autres au prorata. Cela donne $(0.517, 0.083, 0.4)$, ce qui dépasse le plafond de l'actif 1 et ne vérifie pas les conditions.}
\end{solution}

\begin{solution}{iq:iv2:portfolio-construction:23}
Oui~: le rendement espéré total vaut environ 15.5 points de base, contre 10 points de base pour un aller-retour. Le signal de seuil de rentabilité vaut environ 1.3 point de base par jour. Une position détenue un seul jour rapporte 2 et coûte 10, ce qui est une perte.

Le facteur de décroissance est $\rho = 2^{-1/5} = 0.8706$, et la somme des $a\rho^t$ vaut $a/(1 - \rho) = 2/0.1294 = 15.45$.

Le signal de seuil de rentabilité vérifie $a/(1 - \rho) = 10$, donc $a = 1.29$.

Avec une seule journée de détention, il faudrait un signal de 10 points de base par jour.

La durée de détention compte donc autant que la taille du signal. Cela implique aussi une zone de non-transaction~: lorsqu'une position est en place et que le signal ne varie que peu, on ne passe pas d'ordre, car l'information nouvelle ne couvre pas le coût. Un portefeuille doit se rapprocher en partie de sa cible, et le taux d'ajustement dépend de la décroissance et du coût.

\iqwrong{Comparer les 2 points de base du signal aux 5 points de base de coût. La bonne comparaison porte sur ce que le signal rapportera au total, net d'un aller-retour. Et un signal lent peut supporter bien davantage de transactions que ne le suggère son chiffre journalier.}
\iqfollow{Et si un second signal arrivait avant que le premier ait décru~?}{Passer les ordres sur la position nette. Le croisement interne de signaux opposés fait économiser des coûts, et c'est l'un des gains les plus importants d'un portefeuille combiné.}
\end{solution}

\begin{solution}{iq:iv2:portfolio-construction:24}
Le meilleur mélange a un ratio de Sharpe d'environ 1.13, et des parts de risque d'environ 60~\% et 40~\%. Avec une corrélation de 0.8, la seconde stratégie n'apporte rien. Avec une corrélation négative, elle apporte beaucoup.

Pour des stratégies de volatilité unitaire, il faut résoudre $Cx = s$, où $C$ est la matrice de corrélation et $s = (1.0, 0.8)$~: $x = (0.835, 0.549)$, d'où des parts de risque de 60.3~\% et 39.7~\%. Le ratio de Sharpe vaut
\[
\sqrt{s^\top C^{-1}s} = \sqrt{\frac{1 + 0.64 - 2 \times 0.3 \times 0.8}{1 - 0.09}} = 1.129.
\]

La seconde stratégie n'apporte rien lorsque son poids est nul, ce qui se produit pour $\rho = S_2/S_1 = 0.8$. Ses rendements s'expliquent alors par la première, à une pente dont le prix est déjà payé. Notons que ce n'est pas le cas pour $\rho = 1$.

Pour $\rho = -0.3$, le ratio de Sharpe vaut $\sqrt{(1.64 + 0.48)/0.91} = 1.53$. Une corrélation négative le relève donc bien plus qu'un meilleur ratio de Sharpe ne le ferait.

\iqwrong{Allouer un capital égal. L'allocation doit se faire au risque, et les stratégies diffèrent souvent beaucoup en volatilité~: avec un capital égal, une seule d'entre elles mène le jeu.}
\end{solution}

\begin{solution}{iq:iv2:portfolio-construction:25}
Les résultats, avec les poids dans l'ordre, figurent dans le tableau ci-dessous. À 10~\% de volatilité, le portefeuille neutre en bêta a pour poids $(0.226, 0.122, -0.369)$ et rapporte 3.20~\%. En partant de positions nulles, avec les coûts, le portefeuille prend les positions $(0.142, 0.069, -0.225)$.

\begin{center}
\begin{tabular}{lccccc}
\toprule
Portefeuille & Poids & Bêta & Rendement & Volatilité & Sharpe \\
\midrule
Sans contrainte & $(0.244, 0.155, -0.212)$ & 0.205 & 2.77~\% & 7.45~\% & 0.372 \\
Neutre en bêta & $(0.145, 0.078, -0.237)$ & 0 & 2.05~\% & 6.40~\% & 0.320 \\
Neutre en bêta et en dollars & $(0.119, 0.119, -0.237)$ & 0 & 2.02~\% & 6.35~\% & 0.318 \\
\bottomrule
\end{tabular}
\end{center}

\textbf{Sans contrainte.} $w = \Sigma^{-1}\alpha/\lambda$. La corrélation est commune, donc $\Sigma$ est facile à construire, et la résolution donne les poids. Le rendement vaut $\alpha^\top w$, et le bêta $\beta^\top w = 0.205$. Ce portefeuille porte donc un pari de marché non voulu.

\textbf{Neutre en bêta.} On remplace $\alpha$ par $\alpha - g\beta$, avec $g = \beta^\top\Sigma^{-1}\alpha/\beta^\top\Sigma^{-1}\beta = 3.53~\%$. C'est l'alpha diminué de ce qu'explique le bêta. Puis $w = \Sigma^{-1}(\alpha - g\beta)/\lambda$, et $\beta^\top w = 0$.

\textbf{Neutre aussi en dollars.} Deux contraintes, deux multiplicateurs. Avec $G$, matrice formée de $\beta$ et de $\mathbf{1}$, il faut résoudre $(G\Sigma^{-1}G^\top)a = G\Sigma^{-1}\alpha$. Puis $w = \Sigma^{-1}(\alpha - G^\top a)/\lambda$. Le ratio de Sharpe ne bouge presque pas ici, de 0.320 à 0.318~: cette contrainte coûte donc peu. Passer à la neutralité en bêta coûte davantage, de 0.372 à 0.320.

\textbf{Mise à l'échelle.} On multiplie par $10~\%/6.40~\% = 1.562$. Le ratio de Sharpe reste le même.

\textbf{Coûts.} Les trois positions étant conservées et leurs signes connus, un coût absolu abaisse simplement chaque $|\alpha_i|$ de $c$~: $\alpha' = (0.058, 0.028, -0.038)$. On applique donc la même formule neutre en bêta avec $\alpha'$, et l'on obtient $(0.142, 0.069, -0.225)$. Il faut ensuite vérifier qu'aucun signe n'a changé et que chaque position vaut toujours la peine d'être ajustée. L'objectif net vaut environ 0.94~\%. Si une position était assez petite pour que son alpha tombe dans $\pm c$, on ne prendrait aucune position, ce qui correspond à la zone de non-transaction.

\textbf{Vérifications sans solveur.} Les conditions du premier ordre~: $\alpha' - \lambda\Sigma w - g\beta = 0$ pour un certain $g$, pour chaque position. Bêta nul~: $1.2 \times 0.142 + 0.8 \times 0.069 - 1.0 \times 0.225 = 0$, à l'arrondi près. Et les limites~: lorsque $\lambda$ augmente, tout se réduit proportionnellement. Lorsque les coûts augmentent, la plus petite position est la première à disparaître.

Ce que je retiendrais du résultat~: le rendement espéré ne vaut que ce que vaut l'alpha, donc je ferais tourner le calcul avec des alphas rétrécis vers zéro et je montrerais comment le portefeuille évolue. Un portefeuille neutre en bêta avec ces données affiche un ratio de Sharpe espéré d'environ 0.3, ordre de grandeur qu'on ne peut tirer que de trois valeurs. Ce qui compte est l'estimation, et non le chiffre.

\iqwrong{Résoudre le problème sans contrainte, puis couvrir le bêta avec un indice. C'est un autre portefeuille valide, mais ce n'est pas l'optimum, sauf si le résidu de l'indice a les bonnes propriétés. La résolution conjointe utilise les noms eux-mêmes pour neutraliser le bêta.}
\end{solution}
